% Pursuing Stacks (cotes 134-2 to 134-8): archive cote 134-2, site page 71 (PDF page 72).
% appendix (letters to Breen). Draft compared: ps1.tex:1572-1587.
% AI-transcribed by Claude Opus 5.5 (claude-opus-5-5, adaptive thinking, effort high), Message Batches API, batch ps1 (msgbatch_01Rv5hCaNTtUt4JHx1yqoHPL), 2026-11-12; not yet checked by eye.
\scanpage{134-2/71}
\mspage{- \hand{5}6 -}

$X_0$ and $X_{\hand{1}}$ , i.e. of the mere homotopy type of $X_0$ and $X_1$ , except \deleted{for}
\deleted{the fact} that we must make explicit the nature of the left exact functor $\phi$ .
I think that this should be possible, in the context of \emph{schemes} in which I
am placed (technically rather sophisticated), on introducing an ``\'etale
tubular neighbourhood'' of $X_0$ in $X_1$ (which is a very \deleted{satisfying}\inserted{interesting} topos, but
not associated to a \deleted{sheaf)}\inserted{scheme}\hand{)}. But this technical construction is only a
paraphrase of an \deleted{topological intuition} extraordinarily simple\inserted{topological intuition}\hand{,} which I will
make explicit, supposing, to fix the ideas\hand{,} that the base field is $\mathbb{C}$ and so
one may work with locally compact spaces in the usual sense. The topological
idea behind the hypothesis of equi-singularity is that there exists a
\emph{tubular neighbourhood} $T$ of $X_0$ in $X$ retracting onto $X_0$ and such
that the pair $(X_0 , T)$ over $X_0$ should be a locally trivial bundle, i.e.
that $T\backslash X_0$ is locally trivial over $X_0$ . In fact if $\partial T$ is the ``boundary''
of $T$ , which also should be a locally trivial bundle on $X_0$ , then $T$ over
$X$ is the conic bundle (= bundle where fibres are cones) $(\simeq (\partial T \times I) \coprod_{\partial T} X_{\hand{0}}$
\margin{$\longmapsto$} where $I = [0, 1]$ , $\partial T \to \partial T \times I$ is defined by $x \over{\to}{\mapsto} (x, 1)$ , and
$\partial T \to X_0$ is the projection) then $\overset{\circ}{T} = T\backslash X_0 \simeq \partial T \times [0, 1\hand{[}$ is $X_0$-homotopic
to $\partial T$ . If $X_0$ and $X_{\hand{1}}$ are non singular, then so also will be $\overset{\circ}{T}$ and $\partial T$ ,
which are then topologically smooth fibrations on $X_0$ . Moreover, putting
$\widetilde{X}_1 = X_{\hand{1}}\backslash\overset{\circ}{T}$ , the inclusion $\widetilde{X}_1 \to X_{\hand{1}}$ is a homotopy equivalence, and $X$ can
be recovered, up to homeomorphism, from the diagram of spaces
\diagram{
\begin{tikzcd}
(\overset{\circ}{T} = T\backslash X_0 \cong)\,\partial T \arrow[r, "j", "\text{inclusion}"'] \arrow[d, "\text{fibration}"', "p"] & \widetilde{X}_{\hand{1}}(\cong X_{\hand{1}}) \\
X_0 &
\end{tikzcd}
}
as an amalgamated sum. In terms of this diagram of spaces, the above functor
$\phi$ interprets immediately as
\[ \phi(F_{\hand{1}}) \simeq p_* j^*(\widetilde{F}_{\hand{1}}) \]
where $F_{\hand{1}} \to \widetilde{F}_{\hand{1}}$ is the restriction from $X_{\hand{1}}$ to $\widetilde{X}_{\hand{1}}$ (which is an equivalence
of categories for the envisaged (locally constant) sheaves). Giving
$F = (F_0 , F_{\hand{1}} , u : F_0 \to \phi F_{\hand{1}})$ can then also be made explicit as giving

\mspage{\hand{71}}
